\chapter{Teoretická východiska}
\label{chap:theory}

Tato kapitola shrnuje klíčové teoretické koncepty, na nichž je výzkumná část práce postavena. Postupně jsou představeny základy marketingového výzkumu, teorie spotřebitelského chování, problematika vstupu nového produktu na trh, metodologie senzorické analýzy potravin, přístupy k~měření cenové citlivosti a~principy moderní marketingové komunikace.

% ============================================================================
\section{Marketingový výzkum}
\label{sec:theory_mr}

Marketingový výzkum představuje systematický proces sběru, analýzy a~interpretace dat, který slouží jako podklad pro marketingová rozhodnutí (Malhotra, 2019). Podle Kotlera a~Kellera (2016, s.~121) je marketingový výzkum „systematické plánování, sběr, analýza a~reportování dat a~zjištění relevantních pro specifickou marketingovou situaci, které organizace čelí."

\subsection{Typy marketingového výzkumu}

V~odborné literatuře se tradičně rozlišují dva základní přístupy k~marketingovému výzkumu (Kozel, 2006; Foret, 2012):

\begin{description}
  \item[Kvalitativní výzkum] se zaměřuje na porozumění motivacím, postojům a~pocitům spotřebitelů. Využívá metody jako hloubkové rozhovory, skupinové diskuse (focus groups) nebo projektivní techniky. Pracuje s~menšími vzorky respondentů, ale umožňuje hlubší vhled do zkoumaného jevu.
  \item[Kvantitativní výzkum] se orientuje na měření a~kvantifikaci jevů na větších vzorcích respondentů. Využívá strukturované dotazníky, experimenty a~pozorování. Výsledky jsou statisticky zobecnitelné na cílovou populaci.
\end{description}

Současné výzkumné přístupy stále častěji kombinují oba typy v~rámci tzv.\ smíšené metodologie (mixed methods), která umožňuje triangulaci dat a~komplexnější pochopení zkoumaného problému (Creswell \& Creswell, 2018).

\subsection{Proces marketingového výzkumu}

Malhotra (2019) definuje šest hlavních kroků marketingového výzkumu:
\begin{enumerate}
  \item definice problému a~cílů výzkumu,
  \item návrh výzkumného designu,
  \item sběr dat,
  \item příprava a~analýza dat,
  \item interpretace výsledků,
  \item prezentace zjištění a~doporučení.
\end{enumerate}

% ============================================================================
\section{Spotřebitelské chování}
\label{sec:theory_consumer}

Spotřebitelské chování lze definovat jako studium procesů, které jedinec nebo skupina uplatňuje při výběru, nákupu, používání a~likvidaci produktů, služeb, myšlenek nebo zážitků k~uspokojení svých potřeb (Solomon, 2020).

\subsection{Model rozhodovacího procesu spotřebitele}

Klasický model rozhodovacího procesu spotřebitele zahrnuje pět fází (Kotler \& Keller, 2016):
\begin{enumerate}
  \item rozpoznání potřeby,
  \item hledání informací,
  \item hodnocení alternativ,
  \item nákupní rozhodnutí,
  \item ponákupní chování.
\end{enumerate}

V~kontextu prémiových potravinových produktů hrají klíčovou roli faktory jako vnímání kvality, důvěra ve značku a~emocionální aspekty nákupu (Vysekalová, 2011). U~produktů typu mražené ovoce v~čokoládě je rozhodování spotřebitele navíc ovlivněno senzorickou zkušeností, vizuální prezentací a~sociálními doporučeními prostřednictvím sociálních sítí.

\subsection{Segmentace trhu}

Efektivní segmentace trhu je základním předpokladem úspěšného marketingového plánování. Dibb a~Simkin (2010) zdůrazňují, že kvalitní segmentace musí produkovat segmenty, které jsou měřitelné, dostatečně velké, dostupné, diferencovatelné a~akčně využitelné.

Pro účely této práce jsou definovány dva primární segmenty:
\begin{itemize}
  \item \textbf{Generace~Z a~mileniálové} (18--40 let) — orientovaní na sociální sítě, vizuální stránku produktu, prémiovost a~zážitek,
  \item \textbf{Matky s~dětmi} (25--40 let) — kladou důraz na složení produktu, absenci umělých přísad a~hledají zdravější alternativu odměny pro děti.
\end{itemize}

% ============================================================================
\section{Vstup nového produktu na trh}
\label{sec:theory_market_entry}

Úspěšné uvedení nového produktu na trh vyžaduje systematickou přípravu zahrnující analýzu trhu, konkurenčního prostředí a~cílové skupiny (Kotler \& Keller, 2016). Klíčovým předpokladem je identifikace mezery na trhu a~jasné vymezení konkurenční výhody.

\subsection{Positioning a~diferenciace}

Ries a~Trout (2001) definují positioning jako proces zaujmutí specifické pozice v~mysli spotřebitele. Efektivní positioning musí jasně komunikovat, v~čem se produkt odlišuje od konkurence a~jakou hodnotu přináší zákazníkovi.

V~případě značky Berrie je navrhovaný positioning postaven na třech pilířích:
\begin{enumerate}
  \item \textbf{Lokální původ surovin} — české ovoce, případně v~\gls{bio} kvalitě,
  \item \textbf{Prémiová kvalita} — technologie šokového mražení (\gls{iqf}),
  \item \textbf{Cenová dostupnost} — konkurenceschopná cena vůči importovanému Franui.
\end{enumerate}

Trout a~Rivkin (2008) zdůrazňují, že diferenciace je v~současném přesyceném tržním prostředí klíčovým faktorem přežití. Produkt, který není schopen se jasně odlišit, ztrácí v~očích spotřebitele svou relevanci.

% ============================================================================
\section{Senzorická analýza potravin}
\label{sec:theory_sensory}

Senzorická analýza je vědecká disciplína zabývající se hodnocením vlastností produktů prostřednictvím lidských smyslů — chuti, vůně, textury, vzhledu a~zvuku (Lawless \& Heymann, 2010). V~potravinářském průmyslu představuje klíčový nástroj pro vývoj nových produktů a~ověřování spotřebitelské přijatelnosti.

\subsection{Central Location Test}

\gls{clt} je metoda senzorického hodnocení, při které respondenti hodnotí produkty na centrálním místě za kontrolovaných podmínek. Výhodou je standardizace podmínek testování, nevýhodou může být nižší ekologická validita ve srovnání s~hodnocením v~domácím prostředí (Ruiz-Capillas et al., 2021).

V~rámci této práce bude \gls{clt} využit k~hodnocení:
\begin{itemize}
  \item textury produktu v~různých časových intervalech po vytažení z~mrazáku (5, 10 a~15 minut),
  \item chuťové rovnováhy mezi čokoládovou vrstvou a~ovocným jádrem,
  \item celkového senzorického dojmu a~preference mezi produktem Berrie a~konkurenčním Franui.
\end{itemize}

% ============================================================================
\section{Cenová citlivost spotřebitelů}
\label{sec:theory_pricing}

Správné nastavení ceny je jedním z~klíčových faktorů úspěšného uvedení nového produktu na trh. Příliš vysoká cena odrazuje potenciální zákazníky, zatímco příliš nízká cena může vzbuzovat pochybnosti o~kvalitě produktu (Kotler \& Keller, 2016).

\subsection{Price Sensitivity Meter}

Metoda \gls{psm}, známá také jako Van Westendorpova metoda, je analytický nástroj pro stanovení optimální cenové hladiny (Van Westendorp, 1976). Respondenti odpovídají na čtyři otázky:
\begin{enumerate}
  \item Při jaké ceně byste produkt považovali za \emph{tak drahý}, že byste o~jeho koupi vůbec neuvažovali?
  \item Při jaké ceně byste produkt považovali za \emph{tak levný}, že byste pochybovali o~jeho kvalitě?
  \item Jaká cena je podle vás \emph{již vysoká}, ale stále byste o~koupi uvažovali?
  \item Jaká cena je podle vás \emph{výhodná} (dobrý nákup)?
\end{enumerate}

Průsečíky kumulativních distribučních křivek odpovědí na tyto otázky definují rozmezí přijatelných cen, optimální cenový bod a~bod cenové indiference.

% ============================================================================
\section{Marketingová komunikace v~digitální éře}
\label{sec:theory_marcom}

Moderní marketingová komunikace je charakterizována přesunem důrazu z~tradičních médií směrem k~digitálním kanálům (De Pelsmacker et al., 2018). Pro produkty cílící na mladší demografické skupiny hrají klíčovou roli sociální sítě, které umožňují virální šíření obsahu a~přímou interakci se spotřebiteli.

\subsection{Role sociálních sítí}

Vrontis et al.\ (2021) ve své systematické review poukazují na rostoucí význam influencer marketingu, zejména na platformách Instagram a~TikTok. Pro kategorie \gls{fmcg} a~potravinářských produktů je vizuálně atraktivní obsah na sociálních sítích často prvním kontaktem spotřebitele s~produktem.

Úspěch značky Franui na evropském trhu je přímo spojen s~virálním obsahem na TikToku, kde se produktová videa šířila organicky díky vizuální přitažlivosti kombinace zmraženého ovoce a~tající čokolády. Tento případ dokládá sílu digitálního \emph{word-of-mouth} marketingu v~potravinářském segmentu.

\subsection{Integrovaná marketingová komunikace}

Chaffey a~Ellis-Chadwick (2019) zdůrazňují nutnost integrace online a~offline komunikačních kanálů. Pro novou značku vstupující na trh je klíčové:
\begin{itemize}
  \item vybudování povědomí prostřednictvím sociálních sítí,
  \item podpora v~místě prodeje (sampling, ochutnávky),
  \item spolupráce s~relevantními influencery v~oblasti food a~lifestyle,
  \item využití e-commerce kanálů (např.\ Rohlík.cz) jako vstupní distribuce.
\end{itemize}
